\chapter*{Abstract} % no number
\label{abtract}

\addcontentsline{toc}{chapter}{Abstract} % add to index


The additive manufacturing technology (i.e. 3D printing) has gained substantial popularity in recent years,
both in an industrial settings (where technologies such as SLA and SLS dominate)
and in the consumer space (where FDM or Fusing Deposition Modelling printers are 
by far the most common option due to the lowe price).

The fast adoption of 3D printers mostly stem from the fact that
they allow to manufacture very complex shapes that would often be impossible to produce
using more traditional manufacturing techniques; However although they are very flexible
in the variety of producible shapes, 3D printers still can't construct any arbitrary
shape.

The biggest challenge that prevent an object to be printed in one piece
is gravity, and the impact it has on the layer segmentation procedure.
The vast majority of 3D printers works by constructing a particular object layer by layer,
where every layer is constructed by ``adding'' the material the desired places.
This of course create some issues when a large area need to be created
without having anything below (often referred to as ``overhangs''); 
In this case the layer cannot be deposited on top of anything,
meaning that most of the time it will foll down and cause a failed print.

The most common way to deal with overhangs is through the use of supports; 
A support is simply something that is printed together with a specific object,
but is only there to sustain a specific overhang region while it prints, and will
then be removed. The vast majority of support structures are not designed by hand,
but they rather automatically generated by specialized algorithms.

Supports are inherently wasteful, as they are often thrown away the moment that they are
detached from the printed object; This create a requirement on support structure optimization
algorithm, namely the objective of producing a support that minimizes the material combustion.
Some research in better support structure optimization algorithms already exist, both
in form of academic literature, and also as part of many popular open-source 3D printing softwares.

One key thing to understand however, is that support structures are not necessarily portable
across different 3D printing technologies,
Indeed some technologies (such as for example SLA, and to a minor extend SLS) exert minimal
forces to the object while it's being printed, meanwhile in DFM, the nozzle will always
put a bit of force onto the object is printing while it deposits the material. This means
that printing tall and skinny structures on a DFM printer often results in them bending
and ultimately causing the print to fail.
SLA and SLS 3D printing software often find themselves designing supports with
tall and skinny structures, as the thinker a support is, the least material ti consume.
FDM softwares however tend to construct structures that are a bit thicker, as support that are too
thin would result in them bending and or braking during printing.

A promising approach for reducing the material necessary to print a support is the use of 
Genetic Algorithms; Indeed this class of algorithms often excels at complex minimization
problems, where exact solutions are too complex to find. In literature we can find 
some examples of this algorithmic approach being applied for support generation,
however all the attempts were designed to work on SLS and SLA 3D printers,
while FDM is left unaddressed by the current literature.

Our work's aim is to fill the void in the literature by creating an genetic-algorithm
based support structure optimizer that is specifically designed with FDM 3D printers 
in mind. For achieving this objective we began by selecting a ``topology'' of the support 
structure. All of the previous literature had used tree-like structures with a fixed
size beam, this is fine for SLA and SLS but when testing this topology on a FDM printer
we found out that the tree wasn't stiff enough, and resulted in a failed print.
Meanwhile open source FDM 3D printing software is also using tree topologies,
but the thickness of the beams is greater at the bottom, making the entire structure
stable (but consuming more materials).
In our algorithm we decided to use a rigid frame topology (where unlike a tree
andy point can have multiple connections to the ground) with fixed-size beams,
as some early testing has shown promising results both in terms of stiffness and in terms
of material usage.
With the topology selected we then preceded designing an algorithm that was able to
optimize this topology based on a cost function.

One key component of the cost function is stiffness, as we want to prevent beam to bend and cause
a failure. however stiffness evaluation of a complex rigid frame is a time consuming process, and 
four our genetic algorithm to converge we need to run this evaluation process hundred of thousands
of times. For this reasons, instead of using traditional stiffness evaluation algorithms, we designed
our own faster alternative. Our algorithm is a monte-carlo approach, as we had to sacrifice a bit of 
accuracy in order to get the required speed. Our tests have shown that (for our specific use-case)
our algorithm is about 400 times faster, while only losing about 4\% of accuracy when compared
to exact methods.

The genetic algorithm was developed in a three stage step. This was done to reduce
the complexity of the problem, by separating it into three separate problems.
The three problems were: Selection of the contact points,
grouping of different areas together, and optimization of the support structures.
The final result is a reduction of approximately 43\% in the material usage
for supports when compared to the state of the art algorithm for FDM 3D printers.
The results where validated in the real world with real print tests, and all the 
models were printed successfully.

Our current approach have a few downsides including longer printing time
and a slightly worst print quality (which are caused by the design of 
the contact points, and we believe can be addressed in future work).
but is overall a substantial improvement over the current state of the art.
We believe that with a few refinements and ease of use improvement
our algorithm could save a measurable amount of plastic if 
it managed to get a wide adoption.



